% !TeX root = ../../main.tex
\subsubsection{Relationship to the legacy two-group statistic}

BARCS and CB\textsuperscript{2} share a sampling principle but not an
implemented statistic.  CB\textsuperscript{2} estimates two weighted group
proportions, permits separate group dispersions, forms their difference, and
uses a Welch--Satterthwaite reference distribution.  BARCS fits a logit
coefficient with one guide dispersion across the full design and uses residual
$m-q$ degrees of freedom.  These choices differ in finite samples even for a
binary design.

An identity-link two-cell generalized least-squares model can be constructed
to reproduce the already-computed CB\textsuperscript{2} group summaries and
variances.  This is an algebraic restatement of those summaries, because the
working precision matrix is defined from CB\textsuperscript{2}'s own weights.
It is useful as a software check but is not a nesting theorem and supplies no
evidence about BARCS calibration.

For small differences away from the probability boundaries, the delta method
makes a logit contrast approximately proportional to a raw-proportion
contrast.  A stronger equivalence would additionally require the groupwise and
pooled dispersion estimators to converge to a common value as the number of
independent biological libraries increases.  Condition-dependent dispersion
violates that assumption, and deeper sequencing at a fixed number of
biological units does not create the required asymptotic regime.  We therefore
make no formal equivalence claim.  The operational generalization is limited
to the mean model: BARCS can estimate continuous, adjusted, and interaction
coefficients that the two-group CB\textsuperscript{2} interface cannot express.
